% Abstract

Mobile-assisted language learning (MALL) applications are now the most widespread medium of independent language study, yet most of them rest on uniform, decontextualised drills decorated with points and streaks, and first-week attrition is reported to exceed 60\%. This thesis asks whether the media that learners already consume habitually --- short-form video and hypercasual games --- can be turned into an effective, personalised vocabulary-learning environment. It presents \textit{Glosy}, a cross-platform mobile application that currently supports English, Greek, and Farsi in any native/learning combination, with a dictionary of roughly 41,000 words tagged by proficiency level.

The central design element is a word-level learner model: every word a learner tracks carries a Free Spaced Repetition Scheduler (FSRS) memory state (difficulty, stability, and retrievability) that is updated on the device after each graded review, so the application works offline and synchronises when a connection returns. Learners watch reels --- videos of at most 90 seconds with timestamped, tokenised, translated subtitles, which users can also create --- and add words to their vocabulary by tapping them. A three-stage cascade recommender selects the reels. A comprehensibility filter keeps only reels of which the learner already knows at least 98\% of the vocabulary; a spaced-repetition stage ranks reels containing words that are due for review first, so that watching a reel doubles as a graded review; and a content-based stage orders the remainder by the TF-IDF cosine similarity of each reel to a profile built from the learner's engagement. Two hypercasual games, a multilingual Wordle-style game and a procedurally generated crossword called Word of Wonders, draw their words from the same tracked vocabulary, so that each round is also a spaced-repetition review.

No controlled study with learners has yet been conducted, so pedagogical effectiveness is not claimed; instead, the algorithmic components are evaluated directly. Over 4,200 benchmark trials the crossword generator never failed, and the richness of its boards was limited by its letter budget rather than its word cap. \textit{Glosy} also faces the cold-start problem of a new application: with no user base yet, no users have uploaded reels, so its content library is small and the recommender has no interaction history to personalise from. The thesis closes with the main limitations --- the absence of a user study and offline annotation of new reels --- and with directions for future work: automated annotation, cold-start mitigation, and a randomised controlled evaluation.

\textbf{Key words:} mobile-assisted language learning, short-form video, spaced repetition, recommender systems, hypercasual games, vocabulary acquisition
